\begin{papercard}{simon2023}
\cf{Stages} \stagetag{A} \stagetag{W} \stagetag{I} (also \stagetag{S}). A systematic review of insomnia apps in the German and British stores (Ulm MHAD group), with a two-level eligibility screen, dual MARS-G rating and a \emph{therapeutic safety} extension. Relevant for the sleep scope and for functional inclusion rules.

\cf{Aim} An updated overview of the quality (MARS-G) of insomnia apps currently in Google Play and the App Store; secondary aims were content, evidence base and potential therapeutic benefit \pc{simon2023}{2}{Therefore, the primary aim of the study was}. Earlier reviews used single or few terms and may be outdated given market volatility \pc{simon2023}{2}{A more complex search strategy may}.

\cf{Search} A validated web crawler screened the European App Store and Google Play with insomnia terms (Table S1), chosen so that included apps represent what help-seekers encounter \pc{simon2023}{2}{A web crawler (Stach et al., 2020) was used to systematically screen}. Search ran 18--23 September 2020. It was supplemented by apps named in systematic reviews and earlier app evaluations, with all entries pooled in a central database and de-duplicated \pc{simon2023}{3}{ducted from 18 September to 23 September 2020}. The German and British storefronts were searched \pc{simon2023}{5}{The British and German app Stores}.

\cf{Two-level eligibility (Table 1)} Level 1, on store title and description: adults with insomnia symptoms; useful for psychotherapeutic treatment; ``insomnia'' or ``sleep disorder'' in title or description; English, German or French. Excluded: professional-only, relatives-only, other primary disorders, needing another device (e.g.\ smartwatch), tablet-only, blended-care-only. Level 2, after download: at least one CBT-I component \emph{or} a sleep diary/assessment, and fully functional. Excluded: \emph{alarm-clock or sensor-based sleep-tracking only}, music/noise/stories only, eBooks, partial content via other modalities \pc{simon2023}{2}{Eligibility criteria for inclusion}. Psychoeducation-only and diary-only apps were kept as relevant for blended care. Technical problems were verified on a second device (Huawei P10 lite; iPhone 6s) \pc{simon2023}{3}{Eventual technical problems were verified}. 2{,}236 identified, 53 included (37 Android, 16 iOS) \pc{simon2023}{4}{From 2236 identified MHAs, 53 MHAs}.

\cf{Rating procedure} Two independent psychology-trained reviewers per app with MARS-G after standardised online training, at least 15 minutes of use, ratings entered in a dedicated database. Apps on both platforms were rated per platform. Item-level differences $\geq$2 points went to a third author and were resolved by discussion \pc{simon2023}{3}{Each MHA was rated by two independent reviewers with a degree in}. Overall ICC (two-way mixed, agreement) = 0.92; $\omega$ = 0.86--0.97 \pc{simon2023}{3}{The overall quality rating showed an excellent level of}.

\cf{Privacy coding} Descriptive features from the MARS-G classification: privacy policy, imprint, passwords, logins. Only information disclosed in the app, its website or its store description was used \pc{simon2023}{3}{The assessment of privacy and security features occurred on a descriptive}. Results: 75\,\% had a privacy statement, 38\,\% required active consent, 70\,\% informed about data handling, 96\,\% had contact or imprint, 13\,\% had emergency functions \pc{simon2023}{6}{Has a privacy statement}.

\cf{Safety extension (method idea)} Beyond MARS-G therapeutic gain, they coded whether apps gave suspected or definitive diagnoses and advised consulting providers, whether sleep restriction came with \emph{contraindications and adverse effects}, differential-diagnosis information, store disclaimers that the app does not replace treatment, and information on on-site help \pc{simon2023}{4}{We extended the criteria defined by}. Only 2 of 16 sleep-restriction apps explained contraindications or adverse effects \pc{simon2023}{6}{of the 16 MHAs featuring sleep restriction}. The information-quality item was rated against the goal stated in the store description \pc{simon2023}{3}{tion quality was evaluated regarding the goal that was defined in the}.

\cf{Analysis} Means of the two raters; subscale means; categories low $<$2.5, moderate 2.5--4, high $\geq$4. Wilcoxon tests by store and price model \pc{simon2023}{4}{The ratings of the two independent reviews were averaged for all cal}. The evidence-base item was checked by searching Google Scholar, developer sites and reviews.

\cf{Limitations stated} A snapshot of September 2020. A meta-level rating rather than per component. 15 minutes is not in-depth. MARS was not validated for insomnia. Goal-relative rating can inflate simple apps. Only German and British stores. Privacy assessed only descriptively \pc{simon2023}{12}{Given the volatility of the app market}.

\cf{Use for this thesis} Directly informs the \emph{sleep boundary} question: sensor-only sleep trackers and alarm clocks were \emph{excluded} because the review targeted therapy, whereas the thesis targets \emph{tracking}. The same evidence can justify the opposite rule if it is stated explicitly. Reusable: the two-level screen (store-level vs.\ in-app criteria), language rule (EN/DE/FR) with dedicated raters, device verification, discrepancy threshold, and the \emph{risk-communication} items (disclaimers, contraindications, adverse effects). These are close analogues of harm and limitation statements.
\end{papercard}
